% Page budget: 1.55 pages | Owns: augmentation topology, additional states, encoder, and closed-loop objective.
% WRITING GUIDE
% Purpose: Define the final augmentation architecture and explain how it is trained inside the known feedback loop.
% Start here: Current implementation decisions D-129, D-140, D-141, D-180, D-220, D-233, and D-237 in docs/decisions.md.
% Supporting material: Quinten's 18 September outline feedback, the Meeting-audit synthesis, Thesis-writeup/Documentation/TRAINING-DESIGN.md, and docs/writeup figures.
% Mandatory papers: Read Hoekstra's 2025 LFR augmentation paper, the 2026 first-principles augmentation paper, the encoder-initialization paper, and Kessels' Chapter 5 before writing this section.
% Research-plan anchor: Preserve Aspect 2's dynamic parallel augmentation objective; document the departure from open-loop training and earlier state routing.
% Current decisions: Separate physical and additional states, keep the controller outside the model, score the complete training window without burn-in, and use one network for the physical correction and added-state update.
% Choices to justify: Final reported routing, added-state dimension, ANN width and depth, zero initialization, encoder history, closed-loop objective, rollout length, full-window scoring, and validation horizon. Use the configuration of the reported thesis run, not a superseded routing experiment.
% Horizon check: Treat encoder history, scored training window, joint-estimation horizon, and free-run validation horizon separately; relate candidates to relevant stable modes and test sensitivity to slower dynamics and free-integrator accumulation.
% Implementation audit: Read scripts/gantry/gantry_interconnect_dynamic.py, gantry_dynamic/{model,config,controller,data,training}.py, and model_augmentation/fit_systems/{interconnect,closed_loop,pre_encoder}.py.
% Reference check: Verify sources for parallel LFR augmentation, encoder initialization, closed-loop state extension, and any claimed architectural precedent against the original paper or thesis.
% Cautions: Earlier routing plans are superseded; distinguish the truth-only hidden absorber from learned states; closed-loop fit alone does not prove autonomous plant recovery.
% Figures: Absorber schematic, author's updated architecture, and thesis-owned common-reference loops. No residual panel or horizon figure.
% Section is finished when: Every signal has a defined frame and timing, the RK4 boundary is explicit, and the described architecture matches the final code.
%
% SOURCE MAP
% Hoekstra et al. 2025 EJC: dynamic-parallel topology, additional states, encoder-based truncated simulation.
% Hoekstra et al. 2026 Automatica submission: extended LFR formulation, baseline-preserving model/encoder initialization, joint identification context.
% Hoekstra et al. 2026 encoder paper: reconstructability-based W^b initialization; it does not identify W^a or test dynamic augmentation.
% Kessels 2025 thesis: direct closed-loop rollout, controller equations, and reconstruction of the machine controller state for each window.
% Drenth 2025 thesis: LPV-specific augmentation precedent; not the unrestricted MLP structure used here.
% This project: gantry specialization, residual-loop implementation, exact routing, RK4 boundary, and verification.
\section{Dynamic LPV-LFR Augmentation}\label{sec:augmentation}

The baseline of Section~\ref{sec:lfr} describes the rigid-body dynamics of
the gantry, but not every effect of the machine.
Garc\'ia-Herreros et al.\ neglect the resonance of the supporting base, the
vibration of the cross-arm on its flexible plates, the detent forces of the
linear motors, and the variation of friction along the
stroke~\cite[Sec.~2.4]{garcia2013model}, and our baseline \eqref{eq:eom}
also leaves out the Coulomb friction that their model contains.
We therefore augment the baseline with a learned component and estimate the
physical parameters $\theta_{\mathrm{base}}$, the network parameters
$\theta_{\mathrm{aug}}$ and the encoder parameters $\eta$ jointly from
records measured under feedback.
Section~\ref{sec:aug_structure} defines the augmented model,
Section~\ref{sec:aug_closed_loop} closes the known controller around it,
Section~\ref{sec:aug_objective} states the identification criterion with
its encoder, and Section~\ref{sec:aug_init} the initialisation.

\subsection{Augmented model}\label{sec:aug_structure}
An omitted vibration mode requires additional model order.
Refitting $\theta_{\mathrm{base}}$ cannot add this order, and neither can a
static correction of the baseline's state update, because both only change
the dynamics of the existing states~\cite[Sec.~2]{hoekstra2025lfr}.
Additional states can, so we identify the dynamic-parallel model of Hoekstra
et al.~\cite[Table~1]{hoekstra2026lfr},
\begin{equation}
\begin{aligned}
 \tilde x_{k+1}
 &=f_{\mathrm{base}}(\theta_{\mathrm{base}},\tilde x_k,\hat u_k)
   +f_{\mathrm{aug}}(\theta_{\mathrm{aug}},\hat x_k,\hat u_k),\\
 \bar x_{k+1}
 &=g_{\mathrm{aug}}(\theta_{\mathrm{aug}},\hat x_k,\hat u_k),\\
 \hat y_k&=h_{\mathrm{base}}(\tilde x_k)=C_n\tilde x_k ,
\end{aligned}
\label{eq:aug_transition}
\end{equation}
where $\tilde x\in\R^6$ is the physical state of Section~\ref{sec:lfr},
$\bar x\in\R^{n_{x_a}}$ collects the additional states,
$\hat x=\col(\tilde x,\bar x)$, and $\hat u_k$ is the input applied in the
closed loop of Section~\ref{sec:aug_closed_loop}.
As in Section~\ref{sec:lfr}, $u=P\uact$ is the generalised force, and
$y=\qact$ are the measured positions.
The learned functions $f_{\mathrm{aug}}$ and $g_{\mathrm{aug}}$ act in
parallel to the baseline transition $f_{\mathrm{base}}$, which keeps its
physical parameters.

Drenth keeps an augmented model in LPV-LFR form, adding states, latent
variables and scheduling variables to the baseline
LFR~\cite[Sec.~5.2]{drenth2025thesis}.
We instead keep the LFR of Section~\ref{sec:lfr} exact and place an
unrestricted network in parallel to its discretised transition.
That transition is one classical fourth-order Runge-Kutta (RK4) step of the
state equation $\dot{\tilde x}=F(\theta_{\mathrm{base}},\tilde x,u)$ of
\eqref{eq:eom} over one sample period $T_s$, denoted
$F_d(\theta_{\mathrm{base}},\tilde x,u)$, with the input held over the step
as the actuator force is held between samples.
Each RK4 stage evaluates $M(Y)$ at its own intermediate state, so $Y$ moves
within the step and the baseline remains self-scheduled.

Because models with neural networks are estimated better from normalised
data~\cite[Sec.~5.3]{hoekstra2026lfr}, all signals and states are
normalised to zero mean and unit standard deviation, and the RK4 step is
wrapped accordingly,
\begin{equation}
 f_{\mathrm{base}}(\theta_{\mathrm{base}},\tilde x,u)
 =S_x^{-1}\bigl(F_d(\theta_{\mathrm{base}},S_x\tilde x+\mu_x,
   S_u u+\mu_u)-\mu_x\bigr),
 \label{eq:aug_fbase}
\end{equation}
where $S_x$, $S_u$ hold the standard deviations and $\mu_x$, $\mu_u$ the
means of the training records; the output is normalised in the same way.
Unlike~\cite[Sec.~5.3]{hoekstra2026lfr}, which takes the state statistics
from a simulation of the baseline, we take them from the recorded positions
and their finite-difference velocities.
The wrapping changes the coordinates of the baseline, not its dynamics, and
from here on $\tilde x$, $\bar x$, $u$ and $y$ denote normalised quantities.

Both learned functions are outputs of one multilayer perceptron (MLP) with
tanh hidden layers, the single learning function
of~\cite[Eq.~(22d)]{hoekstra2026lfr},
\begin{equation}
 \begin{bmatrix}
  f_{\mathrm{aug}}(\theta_{\mathrm{aug}},\hat x_k,\hat u_k)\\
  g_{\mathrm{aug}}(\theta_{\mathrm{aug}},\hat x_k,\hat u_k)
 \end{bmatrix}
 =\phi_{\mathrm{aug}}(\theta_{\mathrm{aug}},z_k),\qquad
 z_k=\col(\hat x_k,\hat u_k),
 \label{eq:aug_mlp}
\end{equation}
where $\theta_{\mathrm{aug}}$ collects the weights and biases of the MLP.
The first six outputs form $f_{\mathrm{aug}}$ and the last $n_{x_a}$ form
$g_{\mathrm{aug}}$.
The correction $f_{\mathrm{aug}}$ acts on all six physical rows.
Restricting it to the rows of $\Theta$, the only axis with stiffness, would
keep the network out of the free $X$ and $Y$ axes, but would leave no direct
correction of the $X$ and $Y$ dynamics, into which the missing dynamics
couple.
The price is that a correction on these rows integrates without a restoring
force, which Section~\ref{sec:aug_closed_loop} addresses.
Because the network acts beside the RK4 step, not inside it,
$f_{\mathrm{aug}}$ is a discrete-time correction to continuous-time physics.

\begin{figure}[t]
 \centering
 \resizebox{\columnwidth}{!}{\input{figures/tex/augmentation_structure}}
 \caption{Augmented model. The physics step $f_{\mathrm{base}}$ and the
 network $\phi_{\mathrm{aug}}$ act in parallel on $\hat x_k$;
 $\phi_{\mathrm{aug}}$ supplies $f_{\mathrm{aug}}$ and $g_{\mathrm{aug}}$,
 and the output map is not learned ($h_{\mathrm{aug}}=0$). The encoder
 $\psi$ sets $\hat x$ at each window start. During training, $u_k$ is the
 closed-loop input $\hat u_k$.}
 \label{fig:aug_structure}
\end{figure}

Fig.~\ref{fig:aug_structure} shows the interconnection: both state parts
enter $\phi_{\mathrm{aug}}$, but only $\tilde x_k$ reaches the output, since
$C_n$ is the kinematic map $\begin{bmatrix}P^\top & 0\end{bmatrix}$ of
\eqref{eq:Ptransform} in normalised coordinates and the output is not
augmented.
Within one step no learned function feeds the input of another, so the
interconnection graph is acyclic; with $F$ smooth wherever $M(Y)$ is
invertible (Section~\ref{sec:lfr}) and tanh activations, the model is well
posed~\cite[Thm.~9]{hoekstra2026lfr}.
The additional states act on the output only through $f_{\mathrm{aug}}$ and
have no assigned physical meaning.
In particular, they are not the coordinates of the absorber in the simulated
plant of Section~\ref{sec:setup}, since any change of coordinates of
$\bar x$ that the network can absorb leaves the output unchanged.

\subsection{Closed-loop simulation}\label{sec:aug_closed_loop}
The records are measured under feedback, and the model is meant to predict
the machine under feedback, with a known and possibly changed controller.
We therefore fit the closed-loop response, rather than an open-loop
simulation driven by the recorded input as
in~\cite[Eq.~(22)]{hoekstra2026lfr}.
A second reason is that $X$ and $Y$ have no stiffness in
\eqref{eq:damping_stiffness_matrices}: in open loop a velocity error, from
the initial state or the learned correction, integrates into a persistent
position error, whereas inside a stable loop the controller acts on it as on
the machine.
Kessels simulates each truncated window of an augmented model with the known
controller in the loop and propagates gradients through
it~\cite[Eqs.~(5.12), (5.13c), (5.13d)]{kessels2025ai}.
We adopt this and keep the controller outside the learned model, so that it
can be replaced in evaluation, as in his test with changed
gains~\cite[Sec.~5.3.2.3]{kessels2025ai}.
The same feedback also suppresses model error, so a small closed-loop error
alone does not show that the model reproduces the plant without the
controller.

\begin{figure}[t]
 \centering
 \includegraphics[width=\columnwidth]{closed_loop_same_reference.pdf}
 \caption{Recorded plant loop (top) and augmented-model loop (bottom),
 driven by one shared reference $r_k$ and feedforward $u_k^{\mathrm{ff}}$,
 with the same feedback controller. The model output is evaluated
 before the input advances its state to the next sample.}
 \label{fig:aug_closed_loop}
\end{figure}

The recorded loop is driven by a reference $r$ and a feedforward
$u^{\mathrm{ff}}$ through a linear time-invariant controller with state $s$
and realisation $(A_s,B_s,C_s,D_s)$, whose matrices include $P$ and the
normalisation.
The model loop receives the same $r$, $u^{\mathrm{ff}}$ and controller and
differs only in its output $\hat y$ (Fig.~\ref{fig:aug_closed_loop}):
\begin{equation}
\begin{aligned}
 s_{k+1}&=A_s s_k+B_s(r_k-y_k), &
 u_k&=C_s s_k+D_s(r_k-y_k)+u^{\mathrm{ff}}_k,\\
 \hat s_{k+1}&=A_s \hat s_k+B_s(r_k-\hat y_k), &
 \hat u_k&=C_s \hat s_k+D_s(r_k-\hat y_k)+u^{\mathrm{ff}}_k ,
\end{aligned}
\label{eq:aug_direct_controller}
\end{equation}
where $u_k$ and $y_k$ are the recorded input and output, and $\hat s$ is the
state of the model's copy of the controller.
Subtracting the recorded loop from the model loop removes $r$ and
$u^{\mathrm{ff}}$,
\begin{equation}
\begin{aligned}
 x^c_{k+1}&=A_s x^c_k+B_s(y_k-\hat y_k),\qquad x^c_\tau=0,\\
 \hat u_k&=u_k+C_s x^c_k+D_s(y_k-\hat y_k),
\end{aligned}
\label{eq:aug_controller}
\end{equation}
where $x^c=\hat s-s$ is the difference of the controller states and $\tau$
is the start of the simulation window.
The model thus receives the recorded input plus the controller's response to
its own output error, without $r$ and $u^{\mathrm{ff}}$ being known.
This is exact when both loops use the same linear controller, not scheduled
on the model state, with the same feedforward and without saturation; a
controller re-discretised at a lower model rate (Section~\ref{sec:setup})
makes it approximate.

The condition $x^c_\tau=0$ starts the model's controller in the state of the
recorded controller, $\hat s_\tau=s_\tau$.
Kessels reaches the same initial condition by reconstructing the machine's
controller state from output, reference and
controller~\cite[Remark~5.4]{kessels2025ai}; here no reconstruction is
needed, because the recorded input $u$ already contains the recorded
controller's action.
Model error before $\tau$ is therefore not carried into the window.
Because the model has no feedthrough, each sample is evaluated in a fixed
order: $\hat y_k$ from the state, then $\hat u_k$ from
\eqref{eq:aug_controller}, then $\hat x$ and $x^c$ advance, which avoids an
algebraic loop despite the controller's feedthrough $D_s$.

\subsection{Identification criterion}\label{sec:aug_objective}
Each window simulates \eqref{eq:aug_transition} and
\eqref{eq:aug_controller} from an initial state, but only the positions are
measured.
Following Beintema et al.~\cite{beintema2023subnet}, as adapted to model
augmentation in~\cite[Eq.~(22e)]{hoekstra2026lfr}, an encoder estimates this
state from the recorded input and output history,
\begin{equation}
 \hat x_{\tau|\tau}
 =\psi_\eta(\xi_\tau),\qquad
 \xi_\tau=\col\bigl(y_{\tau-n_a:\tau},\,u_{\tau-n_b:\tau}\bigr),
 \label{eq:aug_encoder}
\end{equation}
where $n_a$ and $n_b$ set the history lengths.
Unlike~\cite[Eq.~(22e)]{hoekstra2026lfr}, both histories include the current
sample, as in the reconstructability map from which the encoder is
initialised~\cite[Eq.~(15)]{hoekstra2026encoder}.
The encoder estimates all states, the measured positions included, whereas
Kessels takes the positions from the
measurement~\cite[Remark~5.3]{kessels2025ai}.
The encoder is a linear map with a nonlinear
correction~\cite[Eq.~(8)]{hoekstra2026encoder},
\begin{equation}
 \psi_\eta(\xi_\tau)
 =\begin{bmatrix}W^b\\ W^a\end{bmatrix}\xi_\tau+\tilde\psi(\xi_\tau),
 \label{eq:aug_enc_lin}
\end{equation}
where $W^b$ and $W^a$ give the physical and the additional states, and
$\tilde\psi$ is an MLP with tanh hidden layers.
The parameters $\eta$ collect $W^b$, $W^a$ and the weights and biases of
$\tilde\psi$.

Over window starts $\mathcal T$ and window length $n_f$, the parameters
$\vartheta=(\theta_{\mathrm{base}},\theta_{\mathrm{aug}},\eta)$ minimise the
output error
\begin{equation}
 V(\vartheta)=\frac{1}{|\mathcal T|\,n_f}
 \sum_{\tau\in\mathcal T}\sum_{\ell=0}^{n_f-1}
 \norm{y_{\tau+\ell}-\hat y_{\tau+\ell|\tau}}_2^2
 \label{eq:aug_loss}
\end{equation}
subject to \eqref{eq:aug_transition}, \eqref{eq:aug_controller} and
\eqref{eq:aug_encoder}.
This is the truncated objective of~\cite[Eq.~(22)]{hoekstra2026lfr}, with
each window simulated in closed loop as in~\cite[Eq.~(5.12)]{kessels2025ai}.
The physical parameters, in the coordinates of Section~\ref{sec:params},
the network and the encoder are estimated jointly; the controller is fixed.
As in both cited objectives, every sample is scored from $\ell=0$, so the
encoder's estimate itself is fitted.

Three horizons are kept apart: the encoder history $n_a$, $n_b$, the scored
window $n_f$, and the validation horizon, a closed-loop simulation of a
whole record from one encoder estimate.
A force error of the model at the window start reaches the output through
the loop of model and controller in \eqref{eq:aug_controller}, which starts
from $x^c_\tau=0$.
We therefore bound $n_f$ by the decay time of that loop's impulse response,
not by the slowest open-loop mode of the plant, so that such an error shows
in the output before the window ends.
Slower behaviour, including drift on the free $X$ and $Y$ axes, is then
constrained in training only through the loop and is tested by the
whole-record validation; Section~\ref{sec:setup} gives all three values.

\subsection{Initialisation}\label{sec:aug_init}
For augmented baseline models, a random initialisation of the learning
components can be unstable, whereas an initialisation that reproduces the
baseline uses its prior information~\cite[Sec.~5.4]{hoekstra2026lfr}.
Hoekstra et al.\ therefore require that the augmented model behaves as the
baseline at initialisation~\cite[Eq.~(29)]{hoekstra2026lfr}, and we meet
this by setting the output layer of the network to zero,
\begin{equation}
 W_L=0,\quad b_L=0
 \quad\Longrightarrow\quad
 f_{\mathrm{aug}}=0,\quad g_{\mathrm{aug}}=0 ,
 \label{eq:aug_zero_init}
\end{equation}
where $W_L$ and $b_L$ are the weight and bias of the last layer of
$\phi_{\mathrm{aug}}$.
From the same physical initial state, the augmented model then reproduces
the baseline output, and the encoded additional states $\bar x_{\tau|\tau}$
have no effect until the output layer is trained.
From this start, $f_{\mathrm{aug}}$ can also learn changes that a different
$\theta_{\mathrm{base}}$ would produce, so the split between the two
components is not unique~\cite[Sec.~5.2]{hoekstra2026lfr}.
We do not add the parameter regularisation
of~\cite[Eq.~(26)]{hoekstra2026lfr}, which bounds the drift of
$\theta_{\mathrm{base}}$ but does not remove this non-uniqueness;
Section~\ref{sec:obc} constrains the split by construction instead.

The encoder rows $W^b$ are initialised from the
baseline~\cite[Sec.~3.1]{hoekstra2026encoder}, linearised about rest at
$Y=0$ with the nominal parameters of Appendix~\ref{app:params}, as the
example of~\cite[Sec.~4.2]{hoekstra2026encoder} linearises its baseline
about the origin,
\begin{equation}
 A_c(Y)=\begin{bmatrix}0 & I\\ -M(Y)^{-1}K & -M(Y)^{-1}C\end{bmatrix},\quad
 B_c(Y)=\begin{bmatrix}0\\ M(Y)^{-1}\end{bmatrix},\quad
 C_d=\begin{bmatrix}P^\top & 0\end{bmatrix}.
 \label{eq:aug_lin}
\end{equation}
With $(A_d,B_d)$ the zero-order-hold discretisation of $(A_c(0),B_c(0))$
over $T_s$, matching the held input, $W^b$ is the reconstructability
map~\cite[Eqs.~(16), (17)]{hoekstra2026encoder}
\begin{equation}
 W^b=\begin{bmatrix}
  A_d^nO_n^{+} & \;r_n-A_d^nO_n^{+}T_n
 \end{bmatrix},
 \label{eq:aug_wb}
\end{equation}
where $n=n_a=n_b$, $O_n$ is the observability matrix of $(A_d,C_d)$ over $n$
steps, $T_n$ the Toeplitz matrix of input-to-output responses, $r_n$ the
input-to-state map, and the two column blocks act on the output and input
parts of $\xi_\tau$; all matrices are normalised as the signals.
Because the positions are measured, the linear model is observable, so $O_n$
has full column rank and $O_n^{+}$ is a left inverse.
The rows $W^a$ have no baseline counterpart and are drawn randomly, and
$\tilde\psi$ starts with a zero output layer, so the encoder initially
returns the linear reconstruction of the physical states.
This is only a starting value, since all encoder parameters are trained
jointly with the model.
